%! Constructing a Confidence Interval for a Population Proportion — Unit 6, Lesson 6.2 — Warm-Up
\documentclass[10pt]{article}
\usepackage{apstats-article}
\usepackage{apstats-boxes}
\newcommand{\TallMath}[1]{$\displaystyle #1\rule[-1.4em]{0pt}{3.2em}$}

\begin{document}

\pageheader{Unit 6, Lesson 6.2}{Warm-Up}
\namedateperiod

\vspace{0.15in}

A random sample of 400 high school students was asked whether they have a part-time
job. Suppose the true population proportion is $p = 0.35$.

\begin{enumerate}
  \item Describe the sampling distribution of $\hat p$. Include shape, mean, and standard
        deviation (verify conditions first).

        \begin{itemize}
          \item \textbf{Random:} \blank{6cm}
          \item \textbf{Large Counts:} $np = $ \blank{2.5cm} $\ge 10$? \blank{1cm} \quad $n(1-p) = $ \blank{2.5cm} $\ge 10$? \blank{1cm}
          \item \textbf{Shape:} \blank{6cm}
          \item $\mu_{\hat p} = $ \blank{3cm}
          \item $\sigma_{\hat p} = \sqrt{p(1-p)/n} = $ \blank{5cm}
        \end{itemize}

  \item Find the probability that $\hat p$ is greater than 0.39.
        \vspace{0.05in}

        $z = \dfrac{\hat p - p}{\sigma_{\hat p}} = \dfrac{0.39 - \blank{1.5cm}}{\blank{1.5cm}} = $ \blank{2cm}

        \vspace{0.05in}
        $P(\hat p > 0.39) = P(z > \blank{1.5cm}) = $ \blank{2.5cm}

  \item In Lesson 6.2 you will flip the question: instead of knowing $p$ and finding
        probabilities, you will know $\hat p$ and \emph{estimate} $p$. Write one sentence
        describing why this is a useful change of perspective.

        \writelines{2}
\end{enumerate}

\end{document}
